\cotefile{157-5}
\scanpage{341}
% top centre
\mspage{257}
% lines 1–2
Ainsi, pour définir \inserted{les} $X_{\alpha+1}(j)$, il faut \deleted{définir} tenir compte de toutes ces contraintes.

% line 3
Donc

% lines 4–5
$1^\circ)$ Pour $j^{\unsure{\prime\prime|\sharp}}$ fixé, $X_{\alpha+1}(j'')$ doit \deleted{\unsure{recevoir}}\inserted{être but de} des flèches
% line 6
\[
 X_\alpha(j') \unsure{\xrightarrow{t_\alpha(\lambda)}|\overset{t_\alpha(\lambda)}{\rule{2.8em}{0.4pt}}} X_{\alpha+1}(j''),
\]
% lines 7–10
pour les $\lambda:j'\unsure{\longrightarrow|\mathrel{-}}j''$ dans $\operatorname{Ob}J_{/j''}$ -- y compris pour $\lambda=\mathrm{id}_{j''}$, auquel cas $t_\alpha(\lambda)$\inserted{$:X_\alpha(j'')\longrightarrow X_{\alpha+1}(j'')$} \unsure{jouera|joue} le rôle du morphisme de transition $\tau_\alpha(j'')$.

% lines 11–18
$\uline{2}^\circ)$ Si $\alpha$ \deleted{n'est pas \unsure{ord$^{\mathrm{al}}$} limite, i.e. de} \inserted{est de la forme $\alpha'+2$, \deleted{alors}} \deleted{\unsure{ou|soit}} \inserted{i.e} s'il n'est \deleted{pas}\inserted{ni} \unsure{un ord$^{\mathrm{inal}}$} limite, ni \deleted{\unsure{de} la forme $\unsure{\alpha'+1}$, alors les flèches \unsure{qui|prit} la forme $\unsure{\beta+1}$ on peut mettre $\tau_\alpha(\unsure{j'})$} successeur d'un ordinal limite, alors on \unsure{exige} de plus, pour les morphismes de transition $\tau_\alpha(j\unsure{'|''})$, \deleted{\unsure{resp.} $\tau_\alpha(j'')$} d'égaliser les \inserted{couples de} flèches \unsure{prévues|prises} dans \unsure{ii)|(ii)} et (iii) ci-dessus.

% lines 19–20
Comme \unsure{l'ens|l'une} de ces données du $1^\circ)$, ou de ces contraintes du $\uline{2}^\circ)$,
% archive sheet number: 340
\scanpage{342}
% top centre
\mspage{258}
% lines 1–9
est de cardinal $\leq\pi$, et que la catégorie $\unsure{\mathcal E|\mathscr E|\mathcal C}_j$ est \unsure{grde|gde} devant $\pi$, on voit qu'on peut bel et bien trouver $X_{\alpha+\unsure{1}}(j)$ et des $t_\alpha(\unsure{\lambda|j}):X_\alpha(j')\unsure{\longrightarrow|\mathrel{-}}X_{\alpha+1}(j)$, de façon à satisfaire (si $\alpha=\unsure{\alpha'|\beta'}+2$) aux contraintes \unsure{explicitées} du $2^\circ$. Cela achève la \deleted{\unsure{On doit}}\inserted{description de} la construction de $X_{\alpha+1}(\deleted{\unsure{j}}\unsure{\cdot|j})$, en termes des \deleted{contraintes} $X_\beta(\cdot)$ $(\beta\leq\alpha)$ et des $t_\beta(\lambda)$ $(\beta<\alpha)$.

% lines 10–17
Il faut enfin définir $X_\alpha(\cdot)$ pour $\alpha$ un ordinal limite. Cette fois, on lui demande seulement de recevoir les $X_\beta(\deleted{\unsure{j}}\cdot)$ pour $\beta<\alpha$. Comme l'ens $\unsure{I|J}_\alpha$ de ces $\beta$ est de cardinal $\leq\pi$ et que les $\mathcal E_j$ sont \unsure{grdes|gdes} devant $\pi$, on \unsure{p}\deleted{\unsure{eut}} bel et bien trouver les $X_\alpha(j)$\unsure{.|:}
% archive sheet number: 341
\scanpage{343}
% top centre
\mspage{259}
% line 1
On a donc prouvé l'essentiel du

% lines 2–11
\uline{Lemme \unsure{4|2}.10.7} Sous les conditions de 4.10.3 a) b), et pour un système $X(\cdot)\in\operatorname{Ob}\prod_{j\in\operatorname{Ob}J}\unsure{\mathcal E|\mathscr E|\mathcal C}_j$ donné, on peut trouver un système inductif $\bigl(X_\alpha(\cdot)\bigr)_\alpha$ dans $\prod\mathcal E_j$, indexé par les ordinaux $\alpha$ tels que $\mathrm{card}\,\alpha\leq\pi$, \deleted{\unsure{et pour}}\inserted{et} pour toute flèche $\lambda:j'\unsure{\mathrel{-}|\longrightarrow}j$ dans $J$ \deleted{\unsure{une}} des flèches $t_\alpha(\lambda):X_\alpha(j')\unsure{\mathrel{-}|\longrightarrow}X_\alpha(j)$, de telle façon que l'on ait les conditions suivantes\unsure{:|.}

% line 12
\[
 \uline{1}^{\circ})\quad X_0(\cdot)=X(\cdot)
\]

% line 13
$2^\circ)$ Les conditions (i) (ii) (iii) (p. 256)

% lines 14–17
$3^\circ)$ Si on se donne $F$ dans $\operatorname{Ob}\underline{\unsure{\Gamma|\mathscr T}}(\mathcal E/J)$, et des flèches $\unsure{u_j|u(j)}:F(j)\longrightarrow X(j)\underset{\unsure{\lambda|{}}}{=}X_0(j)$, \unsure{sps|tqs} que pour \unsure{tte} flèche $\lambda$\inserted{$:j'\longrightarrow j$} dans $J$, on a
% bottom, label beside diagram
\diagram{\begin{tikzcd}[ampersand replacement=\&,column sep=1.15em,row sep=3ex]
 {} \arrow[d,phantom,"{\text{\unsure{(iv)|(4)}}}" description] \& F(j') \arrow[r,"{u(j')}"] \arrow[d,"{F(\lambda)}"'] \& X_0(j') \arrow[dr,"{t_0(\lambda)}"] \& {} \& {} \arrow[d,phantom,"{\text{\uline{commutatif}}}" description] \\
 {} \& F(j) \arrow[r,"{u(j)}"] \& X_0(j) \arrow[r,"{\tau_0(j)}"] \& X_1(j\unsure{{}^{\prime\prime}|{}}) \& \cdot\hspace{2.6em}\cdot
\end{tikzcd}}
% archive sheet number: 342
\scanpage{344}
% top centre
\mspage{260}
% lines 1–16
Il reste seulement à prouver ce dernier complément $3^\circ$. C'est une contrainte \inserted{uniquement} sur la construction de $X_1(\cdot)$ et des $t_0(\lambda)$. Supposons donc que l'on ait construit $X_1(\cdot)$ sans tenir compte de cette contrainte. Pour $j$ fixé, il faut \unsure{savoir|donc} ``rendre commutatif\unsure{s|}'' tous ces carrés, i.e. toutes ces doubles-flèches $\bigl(t_0(\lambda)u(j'),\,\tau_0(j)u(j)F(\lambda)\bigr)$ de but $X_1(j)$. Comme l'ens de ces doubles-flèches est de cardinal $\leq\deleted{\unsure{\pi}}\unsure{\mathrm{card}|\mathrm{Card}}\operatorname{Ob}\unsure{\mathcal E|\mathscr E|J}_{/j}\leq\pi$, on peut trouver $X_1(j)\longrightarrow X'_1(j)$ dans $\mathcal E_j$, égalisant toutes ces doubles-flèches. \unsure{donc|dc} quitte à remplacer $X_1(j)$ par $X'_1(j)$, on \unsure{gagne|gagné}.
% archive sheet number: 343
\scanpage{345}
% top centre
\mspage{261}
% lines 1–7
Il est temps enfin de faire appel à l'hypothèse c) de 4.10.3. Soit donc $\alpha$ un ordinal limite, comme stipulé dans cette condition. Comme $\mathrm{card}\,\alpha\leq\pi$, \unsure{à|a} fortiori $\mathrm{card}\,\beta\leq\pi$ si $\beta<\alpha$, on a donc un système inductif
% line 8
\[
 \bigl(X_\beta(\cdot)\bigr)_{\beta\in\unsure{I|J}_\alpha}
\]
% lines 9–11
dans $\prod_{j\in J}\unsure{\mathcal E|\mathscr E|\mathcal C}_j$. Comme les $\varinjlim$ de type $I_\alpha$ \unsure{existent|existant} dans $\mathcal E_j$, on \unsure{peut} \deleted{\unsure{co}} donc considérer
% line 12
\[
 G(j)=\varinjlim_{\beta\in I_\alpha}X_\beta(j)\in\operatorname{Ob}\mathcal E_j
\]

% lines 13–14
\unsure{Si} $\lambda:j'\unsure{\mathrel{-}|\longrightarrow}j$ est une flèche de $J$, on en déduit une $\lambda$-flèche\unsure{s|}
% line 15, above display
\deleted{$G(\lambda):G(j)=\unsure{l}$}
% line 15
\[
 G(\lambda):G(j)=\varinjlim X_\beta(j)\longrightarrow
 G(j')\unsure{\simeq|=}\varinjlim X_{\unsure{\beta+1|\beta}}(j')
\]
% line 16
par $G(\lambda)=\varinjlim t_\beta(\lambda)$. Il faut
% archive sheet number: 344
\scanpage{346}
% top centre
\mspage{262}
% lines 1–2
\deleted{La condition (i) p. 256 \unsure{implique|indique} que $G(\mathrm{id}_j)=\mathrm{id}_{G(j)}$, la}
% lines 3–7
faire attention que les $t_\beta(\lambda)$ ne définissent \unsure{pas} un hom de systèmes \unsure{projectifs}, mais par la condition (ii) p. 256, ils définissent un hom de $\ast$\unsure{Ind|ind}-objets, donc un hom. sur les $\varinjlim$.
% left margin beside lines 3–10
% slanted
\margin{Attention, c'est ici qu'on utilise \unsure{tacitement} le fait que les foncteurs $\lambda_*:\unsure{\mathcal E|\mathscr E|\mathcal C}_{j'}\longrightarrow\mathcal E_j$ commutent aux $\varinjlim$ de type $I_\alpha$, ce qui implique que les \unsure{inclusions} $\mathcal E_j\longrightarrow\mathcal E$ commutent à ce type de $\varinjlim$.}
% lines 7–12
La condition (i) assure que les $G(\lambda)$ \unsure{transforment|transforme} identités en identités, et la condition (iii) que l'on a $G(\mu\lambda)=G(\mu)G(\lambda)$. Donc $G$ apparaît comme une section de $\mathcal E$ sur $J$. On a d'autre part \deleted{\unsure{des}} des flèches
% line 13
\[
 X(j)\overset{\mathrm{def}}{=}X_0(j)\longrightarrow
 G(j)=\lim_{\beta\in\unsure{I|J}_\alpha}X_\beta(j),
\]
% lines 14–17
ce qui prouve $\unsure{F|P}(\unsure{1|A})$. D'autre part, la condition $3^\circ)$ du lemme prouve la validité de $F'(\unsure{1|A})$. Cela achève la dém. de 4.10.3, et par \unsure{C$^{\mathrm{t}}$|C$'$} de \inserted{4.10.1 et de} 4.9.3.
% archive sheet number: 345
\scanpage{347}
% top centre
\mspage{263}
% lines 1–8
\uline{Corollaire} 4.10.8 (du \unsure{Th.|th.} 4.10.1) Soit $N\longrightarrow J$ un foncteur cofibrant, tel que les fibres soient des catégories accessibles (4.8), les foncteurs \unsure{coch.} de base $\lambda_*:N_{j'}\longrightarrow N_j$ \unsure{accessible|accessibles}, enfin \deleted{les} $J$ équivalente à une petite catégorie. Alors la catégorie $M=\underline{\unsure{\Gamma|\mathscr T}}(N/J)$ est accessible.
% lines 8–14
Plus précisément, soit $\pi_0$ un cardinal \inserted{infini} tel que \deleted{l'on ait}\inserted{les} \deleted{$\unsure{\mathrm{card}\,J\leq}$}\inserted{catégories} $N_j$ soient $\pi_0$-accessibles, et \inserted{que} les foncteurs $\lambda_*$ soient $\pi_0$-accessibles, et soit $\pi$ un cardinal \inserted{infini} tel que l'on ait \textcircled{a} $\pi^{\pi_0}=\pi$, \inserted{\textcircled{b} $\forall\lambda:j'\longrightarrow j$ dans $J$, $\lambda_*:\unsure{\mathcal E|\mathscr E|\mathcal C}_j\unsure{\mathrel{-}|\longrightarrow}\mathcal E_{j'}$ applique $(\mathcal E_{\unsure{j|j'}})_{\pi_0}$ dans $(\mathcal E_j)_{\unsure{\pi|\pi_0}}$} \deleted{\unsure{et}}\textcircled{\unsure{c|b}} $\mathrm{card}\,\unsure{\mathrm{réd}|\mathrm{red}}\,J\leq\pi$,
% lines 15–16
Alors \deleted{\illegible{1}} $M$ est $\pi$-accessible, et plus précisément, on a
% line 17
\[
 M\unsure{\approxeq|\simeq}\operatorname{Ind}_{\pi}(M_\pi),
\]
% lines 18–20
où $M_\pi$ est la sous-catégorie \unsure{pleine} de $M$ formée des sections \inserted{$F$ de $N$ sur $J$} telles que l'on ait $F(j)\in(\mathcal E_{\deleted{\unsure{\pi}}\unsure{j|\pi}})_\pi\quad\forall j\in\operatorname{Ob}J$,
% archive sheet number: 346
\scanpage{348}
% top centre
\mspage{264}
% lines 1–2
$(\unsure{\mathcal E|\mathscr E|\mathcal C}_{\unsure{j|\pi}})_\pi$ désignant la sous-catégorie pleine de $\mathcal E_j$ formée des objets $\pi$-accessibles de $\mathcal E_j$.

% lines 3–12
En effet, \unsure{les'|les|l'} hypothèses \unsure{du début|dites|essentielles} impliquent l'existence de $\pi_0$ et $\pi$ satisfaisant les conditions dites. Les conditions a) et b) sur $\pi$, \deleted{$\unsure{\pi}$} plus le fait que les foncteurs $\lambda_*$ soient $\pi_0$-accessibles, impliquent que la sous-catégorie \unsure{str.} pleine $N_\pi$ de $N$ \deleted{formée des objets}\inserted{dont les fibres} sont les $(N_j)_\pi$, est une sous-catégorie cofibrée, et que de plus $N$ est $J$-équivalente à $\operatorname{Ind}_\pi(N_\pi/J)$.

% lines 13–18
\deleted{La condition} D'autre part, les fibres de $N_\pi/J$ sont stables par $\varinjlim$ \unsure{grdes|gdes} devant $\pi_0$ et de cardinal $\leq\pi$, donc la condition ``\unsure{ordinale|ordinaire}'' de 4.10.1. est vérifiée \inserted{(en y faisant $D=N_\pi$)}. \deleted{\unsure{D}} Comme de plus $\mathrm{Card}\,J\leq\pi$, on en conclut \inserted{par 4.10.1} que \deleted{$\unsure{N\simeq|M\simeq}$} $M\approx\operatorname{Ind}_\pi(M_\pi)$.
% archive sheet number: 347
\scanpage{349}
% top centre
\mspage{265}
% lines 1–4
\uline{Corollaire} 4.10.9 : Soit $N\unsure{\mathrel{-}|\longrightarrow}J$ une ``cofibration en paratopos'',$^*$ i.e. un foncteur cofibrant tel que les fibres soient des cat. paratopos, et les foncteurs
% upper left beside opening quotation
% slanted
\margin{\unsure{terme|termes} \unsure{ambigu|ambigus}, cf la double \unsure{explicitation}}
% lines 5–10
\unsure{cochgt|cochgt.} de base \deleted{\unsure{soient}}\inserted{soient des foncteurs ``image \deleted{\unsure{direc}}\unsure{réciproque|inverse}'' \unsure{ou|,} (\unsure{resp.} ``image \unsure{directe|inverse}''),} pour des morphismes de paratopos, i.e qu'ils commutent aux (petites) $\varinjlim$\inserted{(\unsure{resp.|ainsi} aux petites $\varprojlim$)}. Supposons de plus $J$ \unsure{équivalent|équivalente} à une petite catégorie. Alors $\underline{\unsure{\Gamma|\mathscr T}}(N/J)$ est une \inserted{catégorie} paratopos.

% lines 11–17
En effet, c'est une catégorie accessible en vertu de 4.10.8, d'autre part l'hypothèse sur les $\lambda_*$ implique aussitôt que \deleted{\unsure{la}} $\underline{\Gamma}(N/J)$ est stable par petites $\varprojlim$ (\unsure{resp} par petites $\varprojlim$). Donc $\underline{\Gamma}(N/J)$ est une catégorie paratopos en vertu de 4.\unsure{5|8}.9 $b\Longrightarrow a$ (resp. $c\Rightarrow\unsure{o|a|0}$).

% lines 18–20
\uline{Corollaire} 4.10.10. Soit $N\mathrel{-}J$ une catégorie \uline{fibrée} (non cofibrée !), telle que les fibres soient des catégories paratopos et les foncteurs
% archive sheet number: 348
\scanpage{350}
% top centre
\mspage{266}
% lines 1–7
changement de base soient des foncteurs images directes pour des morphismes de paratopos (i.e. des foncteurs accessibles, commutant aux \inserted{petites} $\varprojlim$). Supposons de plus $J$ équivalente à une petite catégorie. Alors $\underline{\unsure{\Gamma|\mathscr T}}(N/J)$ est une catégorie paratopos.

% lines 8–16
En effet, l'hypothèse implique que pour \deleted{\unsure{tout|toute} \unsure{$\lambda\longrightarrow$|de}} flèche $\lambda:j'\unsure{\mathrel{-}|\longrightarrow}j$ dans $J$, le foncteur $\lambda_*:N_{j'}\mathrel{-}N_j$ admet un adjoint à gauche $\lambda^*:N_j\mathrel{-}N_{j'}$, d'où s'ensuit que $N\unsure{\longleftarrow|\longrightarrow}J$ est aussi \unsure{cofibrant|cofibré}, les \deleted{\unsure{morphis}} foncteurs de \unsure{cochgt.} de base \unsure{commutant|commutent} aux petites $\unsure{\varinjlim|\underline{\lim}}$. Donc on peut appliquer 4.\deleted{\unsure{8}}\unsure{10|8}.9, qui donne le résultat voulu.

% lines 17–21
\uline{Question} 4.10.11. Soit $N\unsure{\longleftarrow|\mathrel{-}}J$ un foncteur \uline{fibrant}, tel que les $N_j$ soient des catégories accessibles, et les foncteurs $\lambda^*$ de \unsure{chgt.} de base accessibles, $J$ étant \unsure{équivalente|équivalent} à une petite catégorie. \unsure{Je|J-}
% archive sheet number: 349
\scanpage{351}
% top centre
\mspage{267}
% lines 1–15
conjecture que sous ces conditions, $\underline{\unsure{\Gamma|\mathscr T}}(N/J)\overset{\mathrm{def}}{=}\underline{\mathrm{Hom}}_J(J,N)$ est accessible, et qu'il en est de $\unsure{\overline{\mathrm m}|\overline{\mathrm w}}$ de $\underline{\mathrm{Hom}}\,\underline{\unsure{\mathrm{cart}|\mathrm{cat}}}_J(J,N)$. C'est ce qui est prouvé dans SGA 4 I 9.2\deleted{\unsure{2|7}}\unsure{4|7}. sous des hypothèses légèrement plus fortes (à l'aide de l'ancienne notion, trop forte à mon gré, de filtration cardinale). Je serais étonné que la démonstration (un peu \unsure{longuette}) de loc. cit. ne s'applique pas au cas général\unsure{.|?} Mais on aimerait avoir un énoncé qui \unsure{coiffe} celui-ci et 4.10.8, dans le cas d'un foncteur $N\longrightarrow J$ qui \deleted{n'est}\inserted{ne serait} supposé ni fibrant ni cofibrant, mais \unsure{....|...} (???).

% lines 16–17, ruled heading
\noindent\vrule\hspace{0.5em}\parbox{0.93\linewidth}{4.11. \unsure{Application|Applications} à l'existence des limites dans une catégorie de la forme $\operatorname{Ind}_{\pi_0}(\unsure{C|\mathcal C})$.}

% lines 18–19
\uline{Proposition} 4.11.1. \deleted{Soit $\unsure{\Phi|\underline{\Phi}}$ un ensemble (pas \unsure{néc} petit) de petites catégories. \unsure{On}} Soit
% archive sheet number: 350
\scanpage{352}
% top centre
\mspage{268}
% lines 1–12
$C$ une catégorie, équivalente à une petite catégorie, $\pi$ un cardinal infini, posons $M=\operatorname{Ind}_\pi(C)$. Soit $\unsure{\Phi|\underline{\Phi}}$ un ens. de catégories de cardinal réduit $\leq\pi$. On se propose \deleted{\unsure{une famille}}\inserted{d'examiner} si $M$ est ``stable par $\varinjlim$ de type $\Phi$'' (i.e stable par $\unsure{\varinjlim|\underline{\lim}}$ de type $I$, quel que soit $I\in\Phi$). \uline{Il} \deleted{\unsure{Il|C}} \deleted{\unsure{en} est ainsi si} \inserted{suffit pour cela que} $C$ soit stable par $\varinjlim$\deleted{\unsure{ites}} de type $\Phi$, et que de plus : \deleted{$C$ est \unsure{st}} $\exists$ un ordinal \inserted{limite} $\alpha$ avec $\mathrm{card}\,\alpha\leq\pi$ et tel que $C$ soit stable par $\varinjlim$ de type $I_\alpha$ (p. ex. tel que $I_\alpha\in\Phi$).

% lines 13–20
\deleted{Cette condition} Inversement, si $M$ est stable par $\varinjlim$ de type $\Phi$, et si on suppose de plus $C$ \unsure{karoubienne} (i.e \unsure{stable|stables} \deleted{dans $M$} par facteurs directs) alors $C$ est stable par $\varinjlim$ de type $\Phi$. \deleted{Et le foncteur d'inclusion} Dans l'un et l'autre cas, \deleted{$C$ est}\inserted{le} foncteur d'inclusion $C\unsure{\lhook\joinrel\longrightarrow|\mathrel{-}}M$ commute aux $\varinjlim$ envisagées.
% lower left, diagonal beside conclusion
% slanted
\margin{\unsure{Explicité|Expliciter} \unsure{ds|cf} les corollaires \unsure{du|des} \unsure{type|types} 4.13.20\unsure{,|à|--} 4.13.21 (p. 341\unsure{)|}}

% lines 21–22
\uline{Démonstration} : \deleted{\unsure{Sous}} a) Condition suffisante. On applique 4.\unsure{90|9|10}.3 pour conclure que
% archive sheet number: 351
\scanpage{353}
% top centre
\mspage{269}
% lines 1–8
pour \unsure{tte} catégorie $\deleted{\unsure{I}}\unsure{J|I}$ de card. réduit $\leq\pi$ (en particulier pour $\deleted{I}J\in\unsure{\Phi|\underline{\Phi}}$), tout foncteur $\varphi:\deleted{I}J\longrightarrow M=\operatorname{Ind}_\pi(C)$ admet une représentation \unsure{indicielle|indicible} $\Phi:J\times I\longrightarrow C$, avec $I$ \unsure{grd|gd} devant $\pi$. Ceci, \deleted{suffit} plus la stabilité de $C$ par $\varinjlim$ de type $J$, suffit à établir l'existence de la $\varinjlim$ de $\varphi$.

% lines 9–16
La réciproque provient du fait que si $C$ est karoubienne, alors son image essentielle dans $M$ n'est autre que $M_\pi$, laquelle est stable dans $M$ par \inserted{\unsure{ttes|les}} $\varinjlim$ de cardinal $\leq\pi$ représentables dans $M$. \deleted{\unsure{Elle}} Si $M$ est stable par $\varinjlim$ de type $\Phi$, il en est donc de $\unsure{\overline{\mathrm m}|\overline{\mathrm w}}$ de $M_\pi$, donc de $C$, \unsure{qed.|q.e.d.}

% lines 17–21
\uline{Corollaire} 4.11.2. Pour que $M$ soit stable par $\varinjlim$ filtrantes \inserted{\deleted{\unsure{avec|ou} \unsure{petites} \unsure{pour}} (resp. par petites $\varinjlim$)}, il suffit que $C$ le soit par $\varinjlim$ filtrantes \inserted{(resp. par $\varinjlim$)} \uline{de cardinal} $\leq\pi$, et c'est aussi nécessaire si $C$ est karoubienne.

% line 22
Supposons $C$ stable par $\varinjlim$ filtr. de
% archive sheet number: 352
\scanpage{354}
% top centre
\mspage{270}
% lines 1–5
\inserted{(resp. par $\varinjlim$)} de cardinal $\leq\pi$, donc par 4.11.1 $M$ est stable par ces $\unsure{\overline{\mathrm m}^{\mathrm{es}}|\overline{\mathrm m}\mathrm s}$ \inserted{types de} limites inductives. \deleted{Il s'ensuit} D'autre part il est stable par $\varinjlim$ \unsure{grdes|gdes} devant $\pi$. Il suffit donc de prouver le

% lines 6–11
\uline{Lemme} 4.\unsure{14|11}.3. Soit $M$ une catégorie, $\pi$ un cardinal. Supposons que $M$ soit stable par $\varinjlim$ filtrantes de cardinal $\leq\pi$ \inserted{(resp par \deleted{\unsure{$\varinjlim$ de}} $\varinjlim$ de cardinal $\leq\pi$, resp \unsure{par sommes|sommes} de cardinal $\leq\pi$)}, \uuline{et} par $\varinjlim$ \unsure{grdes|gdes} devant $\pi$. Alors $M$ est stable par \inserted{petites} $\varinjlim$ filtrantes, (resp. par \inserted{petites $\varinjlim$, resp. par} \deleted{petites} petites sommes).

% line 12
Ce \inserted{lemme se} prouve comme 4.4.14 page 97.

% lines 13–14
Reste la réciproque dans 4.\unsure{11|10}.2, qui résulte de l'assertion similaire dans 4.11.1.

% lines 15–19
\uline{Corollaire} 4.11.4. Quand $M$\inserted{$=\operatorname{Ind}_\pi C$} est-il stable par sommes ? Si $\exists$ ordinal $\alpha$ tel que $\mathrm{card}\,\alpha\leq\pi$ et que $C$ soit stable par $\varinjlim$ de type $I_\alpha$, alors il suffit que $C$ soit stable par sommes de cardinal $\leq\pi$. Cette condition est aussi nécessaire \unsure{{}|si} $C$ est karoubienne.
% lower left beside last corollary
% slanted
\margin{Pour \unsure{autres} \unsure{variantes|variations} de 4.11.1, cf 4.13.19 (p. 340) et corollaires.}

% bottom line
Se prouve comme 4.11.2, en utilisant 4.11.3.
% archive sheet number: 353
\scanpage{355}
% top centre
\mspage{271}
% lines 1–4
\uline{Proposition} 4.11.5. Soit $\unsure{\underline{\Phi}|\Phi}$ un ensemble de \deleted{petites} catégories \inserted{$I$} équivalentes à des \deleted{petites} catégories \inserted{finies} et ``rigides'' i.e. pour tt $i\in\operatorname{Ob} I$, \deleted{\unsure{tout}} $\operatorname{Hom}_I(i,i)=\{\mathrm{id}_i\}$.
% upper left beside proposition
% slanted
\margin{? Démonstration pas achevée}
% lines 4–12
\deleted{Soit $C$} \unsure{Soit|Soient} $C$ une catégorie équivalente à une petite catégorie, $\pi$ un cardinal infini, $M=\operatorname{Ind}_\pi C$. Pour que $M$ soit stable par $\varinjlim$ (resp. par $\varprojlim$) de type $\Phi$, il \uline{suffit} que $C$ le soit. Dans le cas \unsure{non respé|non resp}, si $C$ est karoubienne, cette condition est aussi nécessaire.

% lines 13–16
Cette dernière assertion est contenue dans l'assertion similaire de 4.11.1. Pour la première assertion, elle résulte de la façon habituelle du

% lines 17–20
\uline{Lemme} 4.11.6. Soit $J$ une catégorie équivalente à une \deleted{petite} catégorie finie, $J$ rigide, \deleted{Alors} $C$ une catégorie équiv. à une petite catégorie, $\pi$ un cardinal infini.
% lower left beside lemma
% slanted
\margin{? Démonstration \unsure{non|pas} \unsure{explicitée}, c'est peut-être faux !}
% lines 21–24
Alors \uline{tout} foncteur \deleted{\unsure{$I$|$J$}} $J\longrightarrow M\overset{\mathrm{déf}}{=}\operatorname{Ind}_\pi C$ admet une représentation indicielle $J\times I\longrightarrow C$, avec $I$ grde devant $\pi$, en d'autres termes le
% archive sheet number: 354
\scanpage{356}
% top centre
\mspage{272}
% line 1
foncteur canonique pl. fidèle
% line 2
\[
\operatorname{Ind}_\pi(\underline{\operatorname{Hom}}(J,C))\longrightarrow\underline{\operatorname{Hom}}(J,M)
\]
% lines 3–5
déduit de l'inclusion $C\unsure{\mathrel{\lhook\joinrel\relbar}|\hookrightarrow} M$ \deleted{est} d'où $\underline{\operatorname{Hom}}(J,C)\unsure{\mathrel{\lhook\joinrel\relbar}|\subset}\underline{\operatorname{Hom}}(J,M)$ (cf. 4.9) est une équivalence de catégories.

% lines 6–7
(\uline{NB}. J'ignore s'il en est encore ainsi, sans hyp de rigidité sur $J$.)

% lines 8–11
\uline{Démonstration} Cela \deleted{résulte du} \inserted{se prouve comme le} résultat \deleted{\unsure{coroll|coroll.}} similaire pour $\operatorname{Ind}(C)$ \unsure{ald|a l d} $\operatorname{Ind}_\pi C$, \deleted{\unsure{prouvé dans|prouvé ds}} \inserted{cf} SGA 4 I 8.8.5, \deleted{\unsure{et div|et dans}}

% lines 12–23; diagonally cancelled block
\deleted{\uline{Corollaire} 4.11.7. Soient $C,\pi,M$ comme dans 4.11.6, et soit $J$ une catégorie \deleted{telle que $\unsure{\operatorname{card}\,\mathrm{red}\,J|\operatorname{card}\,\mathrm{réd}\,J|\operatorname{card}\operatorname{Fl} J|\operatorname{card} J}\leq\pi$\unsure{)|}} \inserted{équivalente à une catégorie finie,} enfin soit $\varphi\colon J\longrightarrow M$ un foncteur. Pour que $\varphi$ admette une ``$\pi$-représentation indicielle'', i.e. pour que $\varphi$ soit dans l'image \inserted{essentielle} de $\operatorname{Ind}_\pi\underline{\operatorname{Hom}}(J,C)\longrightarrow\underline{\operatorname{Hom}}(J,M)$, il (\unsure{f. et|f. et.}) suffit que $\varphi$ admette une représentation indicielle sans plus, i.e. que son image dans $\underline{\operatorname{Hom}}(J,\operatorname{Ind}(C))$ soit dans l'image essentielle de $\underline{\operatorname{Ind}}(\underline{\operatorname{Hom}}(J,C))$.}
% archive sheet number: 355
